\titledquestion{Dichotomy}

$\sat{3}$ is a well-known $\complete{\NP}$ problem and it has many variants. A kind of dichotomy theorems exists because the complexity problems are mostly either in $\P$ or is $\complete{\NP}$. Schaefer's dichotomy theorem shows that a large set of $\SAT$-like problems are $\complete{\NP}$, while only $6$ kinds of problems can be solved in polynomial time.

First of all we should claim some notations of $\SAT$(the Boolean satisfiability problem). Given a variable set $U = \{u_1,u_2,\cdots,u_n\}$, a \textbf{boolean formula} is defined as the combination of unary/binary operators $\vee$(OR), $\wedge$(AND), $\neg$(NOT) and the variables in the variable set $U$. Given a boolean formula $\phi$, if there exists a variable combination that makes $\phi$ true, then $\phi$ is satisfiable, otherwise $\phi$ is unsatisfiable.

A \textbf{SAT formula} is a logical formula in conjunctive normal form (CNF) Specifically, a \textbf{3-SAT formula} $\phi$ is a conjunction (AND) of clauses $C_1, C_2, \ldots, C_m$, and each clause $C_i$ is a disjunction (OR) of three literals (variables or their negations) from $X \cup \{\neg x_1, \neg x_2, \ldots, \neg x_n\}$. We will give out the initial $\complete{\NP}$ problem: $\sat{3}$.

$\sat{3}$: Given a set of Boolean variables $X = \{x_1, x_2, \ldots, x_n\}$ and a \textbf{3-SAT formula} $\phi$ on $X$, determine whether there exists at least one truth assignment $\tau$ that makes the formula $\phi$ evaluate to true. The yes-instance of $\sat{3}$ is:
\begin{equation*}
\mathsf{3\text{-}SAT} = 
\left\{\langle{\phi\rangle} \middle| 
\begin{aligned}
& \phi \text{ is a 3-SAT formula with variables } X = {x_1, x_2, \ldots, x_n}\\
&\text{and clauses } C_1, C_2, \ldots, C_m \text{ such that there exists} \text{ a truth } \\
& \text{assignment }\tau: X \rightarrow \{\text{true}, \text{false}\}\text{ such that } \tau(\phi) = \text{true}.
\end{aligned}
\right\}
\end{equation*}
In this problem, you are supposed to give out a reduction from the following problem to those problems in $\complete{\NP}$ if it is in $\complete{\NP}$, otherwise give out a polynomial time algorithm to solve it. Moreover, you are required to give out the yes-instances of those problems if they are in $\complete{\NP}$.
\begin{parts}
    \part[5] \textsf{Horn-SAT}: Given a set of Boolean variables $X = \{x_1, x_2, \ldots, x_n\}$ and a \textbf{Horn-formula} $\phi$ on $X$, determine whether there exists at least one truth assignment $\tau$ that makes the formula $\phi$ evaluate to true. A \textbf{Horn-formula} may contain many \textbf{Horn-clause}s.
    A \textbf{Horn-clause} is either a single literal (a positive or negative variable) or a disjunction of at most one positive literal and one or more negative literals. %In other words, Horn clauses are either unit clauses (single literals) or implications of the form $(A \rightarrow B_1 \land B_2 \land \ldots \land B_k)$, where $A$ is a positive literal and $B_1, B_2, \ldots, B_k$ are negative literals.

\begin{solution}
    The Horn-SAT problem is in $\P$, and we can solve it in polynomial time as follow:
    \begin{itemize}
        \item First, search for clauses without negative literals, if there is no such clause, return true because we can set all variables to false to make the formula true.
        \item If there is a clause, let's say $C_i$, without negative literals, because the Horn-formula only allows at most one positive literal in a clause, we can set the variable,let's say $x_i$ in $C_i$ to true and remove $C_i$ from the formula. Then we tranverse all other clauses, let's say $C_j$.
        \item If $C_j$ contains $x_i$, we remove $C_j$. If $C_j$ contains $\neg x_i$, we remove $\neg x_i$ from $C_j$. If $C_j$ contains no $x_i$ or $\neg x_i$, we do nothing.
        \item Repeat the above steps until there is no clause in the formula or there appears an empty clause. If there is no clause in the formula, return true. If there is an empty clause, return false.
        \item Because all the above steps could be done in polynomial time, the Horn-SAT problem is in $\P$.
    \end{itemize}
\end{solution}
\newpage
    \part[5]
    $\sat{4}$: Given a set of Boolean variables $X = \{x_1, x_2, \ldots, x_n\}$ and a \textbf{4-SAT formula}(a conjunction of clauses where each clause $C_i$ is a disjunction of $4$ literals) $\phi$ on $X$, determine whether there exists at least one truth assignment $\tau$ that makes the formula $\phi$ evaluate to true.
\begin{solution}
    The $\sat{4}$ problem is in $\complete{\NP}$, and we can reduce $\sat{3}$ to $\sat{4}$ as follow: Given a 3-SAT formula $\phi$ on $X$, we want to construct a 4-SAT formula $\phi'$ on $X$. For each clause $C_i$ in $\phi$, $C_i$ contains $3$ literals, we add a new variable $y_i$ and a new clause $C_i' = C_i \vee y_i$ to $\phi'$. This could be done in polynomial time, the $\sat{3}$ problem is polynomial-time reducible to $\sat{4}$, and the $\sat{4}$ problem is in $\complete{\NP}$.
\end{solution}
    \part[5]
    $\sat{(3,3)}$: Given a set of Boolean variables $X = \{x_1, x_2, \ldots, x_n\}$ and a \textbf{3-SAT formula} $\phi$ on $X$, where \textbf{each variable appears at most three times}, determine whether there exists at least one truth assignment $\tau$ that makes the formula $\phi$ evaluate to true.\\
    \textbf{Hint:} Consider how $a\Leftrightarrow b$ can be transformed into CNF. 
\begin{solution}
    The $\sat{(3,3)}$ problem is in $\complete{\NP}$. To reduce a general 3-SAT instance $\phi$ to a $(3,3)$-SAT instance $\phi'$ (where each variable appears at most 3 times), do the following:\\
    For each variable $x$ in $\phi$ that appears more than 3 times, suppose x appears in k clauses. Create k new
        variables$\{ x_1 \dots x_k \}$ and replace the ith occurrence of x with $x_i$. Append
        the clause$ \{x \cup \not x_{i+1}\}$ for $i=1,...,k-1$ and the clause $\{x_k \cup \not x_1\}$ to the formula.\\
        In the new instance, the clause $x \cup \not x_{i+1}$
         implies that if $x_i$ is false, $x_{i+1}$ must be
        false as well. The cyclic structure of the clauses therefore forces the xi to be either
        all true or all false, so the new instance is satisfiable if the original one is. Moreover
        the transformation requires polynomial time. 
    In general, this construction takes polynomial time and ensures each variable in $\phi'$ appears at most 3 times in total, giving a valid $(3,3)$-SAT formula.
\end{solution}


\end{parts}